\documentclass[11pt]{article}
\usepackage[a4paper,margin=1in]{geometry}
\usepackage{fontspec}
\usepackage[boldfont=false]{xeCJK}
\setCJKmainfont{FandolSong-Regular.otf}
\usepackage{amsmath}
\usepackage{amssymb}
\usepackage{array}
\usepackage{booktabs}
\usepackage{graphicx}
\usepackage{adjustbox}
\usepackage{enumitem}
\setlist[itemize]{topsep=2pt,partopsep=0pt,parsep=0pt,itemsep=2pt,leftmargin=1.4em}
\setlist[enumerate]{topsep=2pt,partopsep=0pt,parsep=0pt,itemsep=2pt,leftmargin=1.6em}
\usepackage{parskip}
\usepackage{fvextra}
\fvset{breaklines=true,breakanywhere=true,fontsize=\small}
\setlength{\parindent}{0pt}

\begin{document}

\begin{center}
  {\LARGE\bfseries Transport in plants}\\[0.35em]
  {\large Igcse Biology}
\end{center}
\vspace{0.6em}

\section*{Xylem and phloem}

Plants have two transport tissues that run through the roots, stems and leaves:

\begin{itemize}
\item \textbf{xylem} 木质部 — carries water and \textbf{mineral ions} 矿物质离子 up from the roots, and gives the plant \textbf{support} 支撑.

\item \textbf{phloem} 韧皮部 — carries \textbf{sucrose} 蔗糖 and \textbf{amino acids} 氨基酸 to wherever they are needed.

\end{itemize}

In a root, stem or leaf, the xylem and phloem lie side by side in \textbf{vascular bundles} 维管束.

\subsection*{Xylem structure (Supplement)}

Xylem is made of long tubes called \textbf{vessels} 导管, well suited to carrying water:

\begin{itemize}
\item their walls are thick and strengthened with \textbf{lignin} 木质素, which also helps support the plant.

\item the cells are dead and empty (they have no cell contents), so water flows through freely.

\item the cells are joined end to end with no walls across the tube, making one long, continuous pipe.

\end{itemize}

\section*{Water uptake}

Water and mineral ions enter the plant through \textbf{root hair cells} 根毛细胞. The huge number of \textbf{root hairs} 根毛 gives a very large \textbf{surface area} 表面积, which speeds up the uptake of water and mineral ions.

The water then follows this pathway:

root hair cells → root \textbf{cortex} 皮层 cells → xylem → \textbf{mesophyll} 叶肉 cells in the leaf.

You can show this pathway by standing a plant or a white flower in water that contains a coloured \textbf{stain} 染色剂. The stain is carried up the xylem and colours the veins.

\section*{Transpiration}

\textbf{Transpiration} 蒸腾作用 is the loss of \textbf{water vapour} 水蒸气 from the leaves.

\begin{itemize}
\item Water \textbf{evaporates} 蒸发 from the wet surfaces of the mesophyll cells into the \textbf{air spaces} 气腔 inside the leaf.

\item The water vapour then \textbf{diffuses} 扩散 out of the leaf through the \textbf{stomata} 气孔.

\end{itemize}

\subsection*{What changes the rate of transpiration}

\begin{center}
\adjustbox{max width=\linewidth}{%
\begin{tabular}{lll}
\toprule
\textbf{Factor} & \textbf{Transpiration is faster when…} & \textbf{Why} \\
\midrule
\textbf{temperature} 温度 & it is hotter & water evaporates faster \\
\textbf{wind speed} 风速 & it is windier & wind carries the water vapour away \\
\textbf{humidity} 湿度 (Supplement) & the air is drier & a bigger difference makes water diffuse out faster \\
\bottomrule
\end{tabular}}
\end{center}

\subsection*{How water rises up the xylem (Supplement)}

As water vapour leaves the leaf, it pulls more water up behind it. This pull is called the \textbf{transpiration pull} 蒸腾拉力. The water molecules stick to one another by \textbf{forces of attraction} 吸引力, so they are drawn up the xylem together as one long \textbf{column} 水柱 of water.

\subsection*{Wilting (Supplement)}

If a plant loses water faster than it can take it up, its cells become soft and the plant \textbf{wilts} 萎蔫 — the leaves and stem droop. This often happens on a hot, dry, windy day.

\section*{Translocation (Supplement)}

\textbf{Translocation} 转运 is the movement of sucrose and amino acids in the phloem, from \textbf{sources} 源 to \textbf{sinks} 库.

\begin{itemize}
\item sources are the parts that \textbf{release} sucrose or amino acids — for example the leaves, which make sugar by photosynthesis.

\item sinks are the parts that \textbf{use or store} them — for example growing roots, or a storage organ.

\end{itemize}

\section*{Exam tips}

\begin{itemize}
\item Xylem: water and mineral ions, upward only, plus support — its cells are dead, with thick lignin walls. Phloem: sucrose and amino acids, in living cells.

\item Learn the water pathway: root hair cell → cortex → xylem → mesophyll.

\item Transpiration = evaporation from the mesophyll, then diffusion out of the stomata. It is faster when hotter, windier or drier.

\item Water rises by the \textbf{transpiration pull}; the water molecules hold together as a column.

\item Translocation goes from a \textbf{source} (where sugar is made or stored) to a \textbf{sink} (where it is used). The direction can change with the season.

\end{itemize}


\end{document}
